\documentclass[11pt,a4paper]{article}
\usepackage{gtu-el-yazisi}
\begin{document}
\notbilgi{Alan Teorisi}{Yük dağılımları ve Dirac distribüsyonu}{Ders9.pdf}

\sayfa{1}
\textbf{Ders 9}
\section{Değişik Türden yük dağılımları ve Dirac Distribüsyonu}

\begin{center}
\begin{tabular}{ccc}
\underline{hacimsel} & \underline{yüzeysel} & \underline{çizgisel}\\[2pt]
$\rho_v\ [\mathrm{C/m^3}]$ & $\rho_s\ [\mathrm{C/m^2}]$ & $\rho_c\ \left[\frac{\mathrm{C}}{\mathrm{m}}\right]$
\end{tabular}
\end{center}

\elcizimi[0.6]{sekiller/Ders9-s1-1.png}[1]

$x=0$'da $\rho_s(y,z)\ \left[\frac{\mathrm{C}}{\mathrm{m}^2}\right]$ yük yoğunluğu olsun, bu yüzeysel yük yoğunluğuna ilişkin hacimsel yük yoğunluğu
\[
\underset{[\mathrm{C/m^3}]}{\rho_v(x,y,z)}=\underset{[\mathrm{C/m^2}]}{\rho_s(y,z)}\,\underset{[1/\mathrm{m}]}{\delta(x)}
\]

\noindent\textbf{$\delta(x)$'in özellikleri} $\heartsuit$

\noindent 1) $a\le x\le b$, $c\le y\le d$, $e\le z\le f$ prizması içindeki toplam yük $\bigl(\rho_v(x,y,z)=\rho_s(y,z)\,\delta(x)\bigr)$
\begin{align*}
Q&=\int_{x=a}^{b}\int_{y=c}^{d}\int_{z=e}^{f}\rho_v(x,y,z)\,dx\,dy\,dz
=\iiint\limits_{x=a\;y=c\;z=e}^{\;b\quad d\quad f}\rho_s(y,z)\,\delta(x)\,dx\,dy\,dz\\
&=\int_{x=a}^{b}\delta(x)\,dx\int_{y=c}^{d}\int_{z=e}^{f}\rho_s(y,z)\,dy\,dz\ \Rightarrow\
\begin{aligned}[t]
&a,b<0,\ Q=0\\
&a,b>0,\ Q=0\\
&a<0<b,\ Q=\int_{y=c}^{d}\int_{z=e}^{f}\rho_s(y,z)\,dy\,dz
\end{aligned}
\end{align*}
\[
\int_{x=a}^{b}\delta(x)\,dx=\begin{cases}0, & ab>0\\ 1, & a<x<b\end{cases}
\]

\noindent 2) $\rho_v(x,y,z)=f(x)\,\rho_s(y,z)\,\delta(x)$ olsun\\
$f(0)=0$ olursa\\
$\delta(x)$ nedeniyle $x\neq 0$ $\rho_v=0$

\sayfa{2}
\noindent $f(x)$ nedeniyle $x=0$ $\rho_v=0$\\
$f(x)\,\delta(x)=0$ ve $f(0)=0$\\
$f(x)=g(x)-g(0)$ olsun\\
$\bigl[g(x)-g(0)\bigr]\,\delta(x)=0$

\formulkutu{g(x)\,\delta(x)=g(0)\,\delta(x)}

\[
\text{3)}\quad\int_a^b\delta(x)\,dx=\begin{cases}0, & ab>0\\ 1, & a<x<b\end{cases}
\]
$a\to-\infty,\ b\to x$
\[
\int_{-\infty}^{x}\delta(x)\,dx=\begin{cases}0, & x<0\\ 1, & x>0\end{cases}=u(x)\quad\text{birim basamak fonksiyonu}
\]

\formulkutu{\delta(x)=\frac{\partial u(x)}{\partial x}}

\noindent 4)

\elcizimi[0.65]{sekiller/Ders9-s2-1.png}[2]

\[
f(x)=\begin{cases}f_-(x), & x<0\\ f_+(x), & x>0\end{cases}
\]
$x=0$'da süreksizlik var
\[
f(x)=f_+(x)\,u(x)+f_-(x)\,u(-x)
\]
\begin{align*}
f'(x)&=f_+'(x)\,u(x)+\underbrace{f_+(x)\,\delta(x)}_{f_+(0)\,\delta(x)}+f'(x)\,u(-x)\footnotemark+f_-'(x)\,u(-x)\\
&\quad-\underbrace{f_-(x)\,\delta(x)}_{f_-(0)\,\delta(x)}
\end{align*}
\footnotetext{El yazısında satırın ortasındaki bu terim büyük bir çarpı ile çizilerek iptal edilmiş.}

\sayfa{3}
\[
f'(x)=\underbrace{\Big[f_-'(x)\,u(-x)+f_+'(x)\,u(x)\Big]}_{\{f'(x)\}}+\underbrace{\Big[f_+(0)-f_-(0)\Big]\delta(x)}_{[|f|]\,\delta(x)}
\]
\[
f'(x)=\underbrace{\{f'(x)\}}_{\text{Regüler kısım}}+\underbrace{[|f|]\,\delta(x)}_{\text{Tekil (singüler kısım)}}
\]

\[
\operatorname{grad}V=\frac{\partial V}{\partial x}\,\vec e_x+\frac{\partial V}{\partial y}\,\vec e_y+\frac{\partial V}{\partial z}\,\vec e_z\ \Big|,
\]
\noindent $x=0$'da bir süreksizlik olsun

\elcizimi[0.35]{sekiller/Ders9-s3-1.png}[3]

\begin{align*}
\frac{\partial V}{\partial x}&=\left\{\frac{\partial V}{\partial x}\right\}+[|V|]\,\delta(x)\quad+[|V\vec e_x|]\,\delta(x)\\
\frac{\partial V}{\partial y}&=\left\{\frac{\partial V}{\partial y}\right\}\\
\frac{\partial V}{\partial z}&=\left\{\frac{\partial V}{\partial z}\right\}
\end{align*}

\noindent $V$'nin süreksizliği herhangi bir $S$ yüzeyi üzerinde olsaydı,

\elcizimi[0.3]{sekiller/Ders9-s3-2.png}[3]

\[
\operatorname{grad}V=\{\operatorname{grad}V\}+[|V\hat n|]\,\delta(s)
\]

\noindent 5) $x=0$'da süreksizliği olan $\vec E(x,y,z)$ için
\[
\operatorname{div}\vec E=\frac{\partial E_x}{\partial x}+\frac{\partial E_y}{\partial y}+\frac{\partial E_z}{\partial z}
\]
\begin{align*}
\frac{\partial E_x}{\partial x}&=\left\{\frac{\partial E_x}{\partial x}\right\}+[|E_x|]\,\delta(x)\\
\frac{\partial E_y}{\partial y}&=\left\{\frac{\partial E_y}{\partial y}\right\}
\end{align*}

\sayfa{4}
\[
+\ \frac{\partial E_z}{\partial z}=\left\{\frac{\partial E_z}{\partial z}\right\}
\]
\noindent\rule{0.85\linewidth}{0.4pt}
\[
\operatorname{div}\vec E=\{\operatorname{div}\vec E\}+[|E_x|]\,\delta(x)
\]
\[
=\{\operatorname{div}\vec E\}+[|\vec E\cdot\vec e_x|]\,\delta(x)
\]

\noindent Genel durumda bir $s$ yüzeyinde süreksizlik varsa
\[
\operatorname{div}\vec E=\{\operatorname{div}\vec E\}+[|\hat n\cdot\vec E|]\,\delta(s)
\]
\[
\operatorname{rot}\vec E=\{\operatorname{rot}\vec E\}+[|\hat n\times\vec E|]\,\delta(s)
\]

\noindent 6) Süreksizlik $x=0$ yerine $x=h$'da olursa\\
$\delta(x)\to\delta(x-h)$
\[
\int_a^b\delta(x-h)\,dx=\begin{cases}1, & a<h<b\\ 0, & \text{diğer}\end{cases}
\]
\[
f(x)\,\delta(x-h)=f(h)\,\delta(x-h)
\]
\[
\frac{\partial u(x-h)}{\partial x}=\delta(x-h)
\]

\elcizimi[0.55]{sekiller/Ders9-s4-1.png}[4]

\[
f'(x)=\{f'(x)\}+\Big[\underbrace{|f(h+0)-f(h-0)|}_{\Delta f}\Big]\,\delta(x-h)
\]

\noindent 7) Bir $C$ eğrisi $S_1$ ve $S_2$ yüzeylerin ara kesiti olsun. $C$ üzerinde bulunan $\rho_c\ [\mathrm{C/m}]$ için hacimsel yük yoğunluğu
\[
\rho_v(x,y,z)=\rho_c\,\delta(S_1)\,\delta(S_2)
\]
\noindent $S_1: x=a$\\
$S_2: y=b$
\[
\underbrace{\delta(x-a)}_{S_1}\,\underbrace{\delta(y-b)}_{S_2}
\]

\sayfa{5}
\noindent Örneğin $z$ eksenine yerleştirilmiş çizgisel
\[
\rho_v(x,y,z)=\rho_c\,\delta(x)\,\delta(y)
\]

\noindent 8) $(x_0,y_0,z_0)$ noktasında bulunan $q$ noktasal yük için hacimsel yük yoğunluğu
\[
\rho_v=q\,\delta(x-x_0)\,\delta(y-y_0)\,\delta(z-z_0)
\]

\elcizimi[0.45]{sekiller/Ders9-s5-1.png}[5]

\begin{ornek}[1]
$\Delta\!\left(\frac{1}{r'}\right)=-4\pi\,\delta(x-x_0)\,\delta(y-y_0)\,\delta(z-z_0)$ olduğunu gösteriniz
\[
r'=\sqrt{(x-x_0)^2+(y-y_0)^2+(z-z_0)^2}
\]
\end{ornek}

\elcizimi[0.55]{sekiller/Ders9-s5-2.png}[5]

\[
V=\frac{1}{4\pi\varepsilon_0}\,\frac{q}{r'}
\]
\[
\Delta V=-\frac{\rho_v}{\varepsilon_0}\quad\text{Poisson}\qquad\text{hacim içinde yük varsa}
\]
\[
\Delta V=0\quad\text{Laplace denklemi}\qquad\text{hacimde yük yoksa}
\]
\[
\rho_v=q\,\delta(x-x_0)\,\delta(y-y_0)\,\delta(z-z_0)
\]
\[
\Delta\!\left(\frac{1}{4\pi\varepsilon_0}\,\frac{q}{r'}\right)=-\frac{q}{\varepsilon_0}\,\delta(x-x_0)\,\delta(y-y_0)\,\delta(z-z_0)
\]
\[
\frac{1}{4\pi}\,\Delta\!\left(\frac{1}{r'}\right)=-\delta(x-x_0)\,\delta(y-y_0)\,\delta(z-z_0)
\]

\sayfa{6}
\[
\Delta\!\left(\frac{1}{r'}\right)=-4\pi\,\delta(x-x_0)\,\delta(y-y_0)\,\delta(z-z_0)
\]
\noindent\textit{silindirik koordinatlardaki uzunluk} $\downarrow$\\
$\rho^2=(x-x_0)^2+(y-y_0)^2$ olmak üzere
\[
\Delta\!\left(\ln\left(\frac{1}{\rho}\right)\right)=-2\pi\,\delta(x-x_0)\,\delta(y-y_0)
\]

\elcizimi[0.4]{sekiller/Ders9-s6-1.png}[6]

\[
Q=\ell\cdot\rho_c\quad(\ell\ \text{uzunluğu boyunca})
\]
\[
\oint_S\varepsilon_0\,\vec E\cdot d\vec s=\ell\rho_c
\qquad
\begin{aligned}
\vec E&=E_\rho(\rho)\,\vec e_\rho\\
d\vec s&=\rho\,d\phi\,dz\,\vec e_\rho
\end{aligned}
\]
\[
\int_{z=0}^{\ell}\int_{\phi=0}^{2\pi}\varepsilon_0\,E_\rho(\rho)\,\rho\,d\phi\,dz=\ell\rho_c
\]
\[
2\pi\ell\varepsilon_0\,\rho\,E_\rho(\rho)=\ell\rho_c
\]
\[
E_\rho(\rho)=\frac{\rho_c}{2\pi\varepsilon_0\rho}\ \Rightarrow\ \vec E=\frac{\rho_c}{2\pi\varepsilon_0\rho}\,\vec e_\rho
\]
\[
\vec E=-\operatorname{grad}V\ \Rightarrow\ \frac{\rho_c}{2\pi\varepsilon_0\rho}=-\frac{\partial V}{\partial\rho}
\]
\[
\Rightarrow\ V=\frac{\rho_c}{2\pi\varepsilon_0}\ln\left(\frac{1}{\rho}\right)+V_0
\]
\[
\Delta V=-\frac{\rho_v}{\varepsilon_0}\qquad\text{ve}\qquad\rho_v=\rho_c\,\delta(x-x_0)\,\delta(y-y_0)
\]
\[
\Delta\!\left(\frac{\rho_c}{2\pi\varepsilon_0}\ln\left(\frac{1}{\rho}\right)+V_0\right)=-\frac{\rho_c}{\varepsilon_0}\,\delta(x-x_0)\,\delta(y-y_0)
\]
\[
\frac{\rho_c}{2\pi\varepsilon_0}\,\Delta\!\left(\ln\left(\frac{1}{\rho}\right)\right)=-\frac{\rho_c}{\varepsilon_0}
\]
\[
\Delta\!\left(\ln\left(\frac{1}{\rho}\right)\right)=-2\pi\,\delta(x-x_0)\,\delta(y-y_0)
\]

\end{document}
